\section{Biểu diễn Thuật toán và Mã nguồn}
\label{sec:algo_code}

\subsection{Trình bày Thuật toán dạng Mã giả (Algorithm2e)}
Thuật toán~\ref{algo:greedy_iap} minh họa việc sử dụng gói \texttt{algorithm2e} với các từ khóa đã được Việt hóa tự động (\textbf{Đầu vào}, \textbf{Đầu ra}, \textbf{Trả về}):

\begin{algorithm}[H]
\caption{Thuật toán gán ca thi Heuristic cơ bản}
\label{algo:greedy_iap}
\LinesNumbered
\KwIn{Tập hợp cán bộ $I$, danh sách ca thi $J$ đã sắp xếp theo thời gian bắt đầu}
\KwOut{Lịch phân công hợp lệ $M \subseteq I \times J$}

$M \leftarrow \emptyset$\;
\ForEach{$j \in J$}{
    $k \leftarrow 0$\;
    \ForEach{$i \in I$}{
        \If{$i$ không trùng lịch với ca $j$ \textbf{và} $k < \text{Capacity}(j)$}{
            $M \leftarrow M \cup \{(i, j)\}$\;
            $k \leftarrow k + 1$\;
        }
    }
    \If{$k < \text{Capacity}(j)$}{
        \KwRet{\text{Không tìm thấy nghiệm khả thi}}\;
    }
}
\KwRet{$M$}\;
\end{algorithm}

\subsection{Khối Mã nguồn (Code Listings)}
Khối mã nguồn sử dụng theme \textbf{Academic / GitHub Light} đã tinh chỉnh: viền mảnh tinh tế, số dòng nằm ngay ngắn trong khung, phông chữ monospace rõ ràng và màu sắc tương phản trang nhã.

\subsubsection{Ví dụ Mã nguồn Python}
Mã nguồn~\ref{lst:python_solver} minh họa hàm xây dựng ràng buộc và tính điểm mệt mỏi trong Python:

\begin{lstlisting}[language=Python, caption={Hàm kiểm tra trùng lịch và tính phạt mệt mỏi trong Python}, label={lst:python_solver}]
from typing import List, Tuple

def evaluate_schedule(assignments: List[Tuple[int, int]], soft_weights: List[float]) -> float:
    """
    Tinh toan tong chi phi phat dua tren bo trong so mem (*@\textit{(soft weights)}@*).
    """
    w_dist, w_consec, w_fair = soft_weights
    total_penalty = 0.0

    # Kiem tra ca thi lien tiep
    consecutive_shifts = count_consecutive(assignments)
    total_penalty += consecutive_shifts * w_consec

    # Kiem tra chenh lech tai cong viec (*@$L_1$ norm@*)
    imbalance = calculate_workload_variance(assignments)
    total_penalty += imbalance * w_fair

    return round(total_penalty, 4)
\end{lstlisting}

\subsubsection{Ví dụ Mã nguồn C/C++}
Bạn cũng có thể chỉ định bất kỳ ngôn ngữ nào khác như C (Mã nguồn~\ref{lst:c_matrix}):

\begin{lstlisting}[language=C, caption={Khởi tạo ma trận chi phí trong ngôn ngữ C}, label={lst:c_matrix}]
#include <stdio.h>
#include <stdlib.h>

#define MAX_INVIGILATORS 100
#define MAX_SESSIONS 50

int main(void) {
    int cost_matrix[MAX_INVIGILATORS][MAX_SESSIONS];
    
    // Khoi tao ma tran trong so mac dinh
    for (int i = 0; i < MAX_INVIGILATORS; ++i) {
        for (int j = 0; j < MAX_SESSIONS; ++j) {
            cost_matrix[i][j] = (i == j) ? 0 : 1;
        }
    }
    printf("Ma tran chi phi da khoi tao thanh cong.\n");
    return EXIT_SUCCESS;
}
\end{lstlisting}
